\section{Extremal graphs without a clique partition for every odd \texorpdfstring{$\alpha$}{alpha}}\label{sec:construction}

\subsection{Two building blocks}

The \nsone\ enumerated extremal graphs with $s(G)=1$ share one shape. Take a partition attaining $p(G)=\alpha-1$
(Lemma~\ref{lem:p}). In every one of them, $\alpha-2$ of its parts are cliques of orders $3$ and $4$, and the last part
$X$ has independence number $2$. Moreover $X$ is one of only two graphs:
\begin{itemize}
\item $X_6$ (\texttt{graph6} \texttt{ETxO}): six vertices and $8$ edges, two more than $K_3+K_3$;
\item $X_7$ (\texttt{graph6} \texttt{FQjRo}): seven vertices and $11$ edges, two more than $K_3+K_4$.
\end{itemize}
Both are complements of maximal triangle-free graphs. Every part sends exactly three edges to the others, from three
distinct vertices, so the edges between parts form a matching, $(3\alpha-3)/2$ edges in all, and the quotient is a
cubic multigraph on the $\alpha-1$ parts. $X_6$ attaches at one triple of its vertices up to automorphism, and $X_7$ at
three (Figure~\ref{fig:blocks}).

The count explains why these graphs are extremal. $X$ spends two more edges than the two cliques it replaces, and the
quotient on $\alpha-1$ parts needs $(3\alpha-3)/2$ edges instead of $\ceil{3\alpha/2}=(3\alpha+1)/2$. That is two fewer,
so the total is again
\begin{equation}\label{eq:count}
t(n,\alpha)+2+\frac{3\alpha-3}{2}=t(n,\alpha)+\ceil{\frac{3\alpha}{2}}.
\end{equation}
The \ntenthreeb\ graphs with $s(G)=2$ at $(10,3,3)$ have a connected $\alpha$-core, the case (b) of
Proposition~\ref{prop:where}. They are $3$-connected graphs on ten vertices with independence number $3$ and $17$
edges, and none of their spanning subgraphs with the same independence number splits into components.

\begin{figure}[t]
\centering
\begin{tikzpicture}[scale=1.0]
\draw[gray!70] (0.000,1.300) -- (-1.126,-0.650);
\draw[gray!70] (0.000,1.300) -- (-0.000,-1.300);
\draw[gray!70] (0.000,1.300) -- (1.126,-0.650);
\draw[gray!70] (-1.126,0.650) -- (1.126,-0.650);
\draw[gray!70] (-1.126,0.650) -- (1.126,0.650);
\draw[gray!70] (-1.126,-0.650) -- (-0.000,-1.300);
\draw[gray!70] (-1.126,-0.650) -- (1.126,-0.650);
\draw[gray!70] (-0.000,-1.300) -- (1.126,0.650);
\node[circle,fill=linknavy,inner sep=1.8pt,label={[font=\scriptsize]90:0}] at (0.000,1.300) {};
\draw[linknavy,thick] (0.000,1.300) -- (0.000,1.600);
\node[circle,fill=linknavy,inner sep=1.8pt,label={[font=\scriptsize]150:1}] at (-1.126,0.650) {};
\draw[linknavy,thick] (-1.126,0.650) -- (-1.386,0.800);
\node[circle,fill=gray,inner sep=1.8pt,label={[font=\scriptsize]210:2}] at (-1.126,-0.650) {};
\node[circle,fill=white,draw=black,inner sep=1.8pt,label={[font=\scriptsize]270:3}] at (-0.000,-1.300) {};
\node[circle,fill=white,draw=black,inner sep=1.8pt,label={[font=\scriptsize]330:4}] at (1.126,-0.650) {};
\node[circle,fill=linknavy,inner sep=1.8pt,label={[font=\scriptsize]390:5}] at (1.126,0.650) {};
\draw[linknavy,thick] (1.126,0.650) -- (1.386,0.800);
\end{tikzpicture}\qquad\qquad
\begin{tikzpicture}[scale=1.0]
\draw[gray!70] (0.000,1.300) -- (-1.267,-0.289);
\draw[gray!70] (0.000,1.300) -- (0.564,-1.171);
\draw[gray!70] (0.000,1.300) -- (1.267,-0.289);
\draw[gray!70] (-1.016,0.811) -- (-0.564,-1.171);
\draw[gray!70] (-1.016,0.811) -- (1.267,-0.289);
\draw[gray!70] (-1.016,0.811) -- (1.016,0.811);
\draw[gray!70] (-1.267,-0.289) -- (0.564,-1.171);
\draw[gray!70] (-1.267,-0.289) -- (1.016,0.811);
\draw[gray!70] (-0.564,-1.171) -- (1.267,-0.289);
\draw[gray!70] (-0.564,-1.171) -- (1.016,0.811);
\draw[gray!70] (0.564,-1.171) -- (1.016,0.811);
\node[circle,fill=linknavy,inner sep=1.8pt,label={[font=\scriptsize]90:0}] at (0.000,1.300) {};
\draw[linknavy,thick] (0.000,1.300) -- (0.000,1.600);
\node[circle,fill=linknavy,inner sep=1.8pt,label={[font=\scriptsize]141:1}] at (-1.016,0.811) {};
\draw[linknavy,thick] (-1.016,0.811) -- (-1.251,0.998);
\node[circle,fill=linknavy,inner sep=1.8pt,label={[font=\scriptsize]193:2}] at (-1.267,-0.289) {};
\draw[linknavy,thick] (-1.267,-0.289) -- (-1.560,-0.356);
\node[circle,fill=white,draw=black,inner sep=1.8pt,label={[font=\scriptsize]244:3}] at (-0.564,-1.171) {};
\node[circle,fill=white,draw=black,inner sep=1.8pt,label={[font=\scriptsize]296:4}] at (0.564,-1.171) {};
\node[circle,fill=gray,inner sep=1.8pt,label={[font=\scriptsize]347:5}] at (1.267,-0.289) {};
\node[circle,fill=gray,inner sep=1.8pt,label={[font=\scriptsize]399:6}] at (1.016,0.811) {};
\end{tikzpicture}

\caption{The two building blocks, $X_6$ (left) and $X_7$ (right), each the complement of a maximal triangle-free graph,
with independence number $2$ and clique cover number $3$. Navy vertices with a stub: an attaching triple (for
$X_7$, one of three up to automorphism, the others being $\{1,2,3\}$ and $\{1,2,6\}$). Open vertices $3,4$ of $X_6$ and
$3,4$ of $X_7$: an independent pair with no neighbour outside the block. Grey vertices: the remaining vertices.}
\label{fig:blocks}
\end{figure}

\subsection{The construction}

Let $\alpha\ge3$ be odd and $3\alpha<n\le4\alpha-1$, and write $n=3\alpha+r$ with $1\le r\le\alpha-1$; $T(n,\alpha)$
has $r$ cliques of order $4$ and $\alpha-r$ of order $3$. Replace two cliques of order $3$ by $X_6$ (possible when
$\alpha-r\ge2$), or a clique of order $3$ and one of order $4$ by $X_7$ (possible when $r\ge1$); since $r\ge1$, at least
one replacement is always available, and the data use both. Take a loopless cubic multigraph $Q$ on the
resulting $\alpha-1$ parts that is $3$-edge-connected. Such a $Q$ exists because $\alpha-1$ is even: the triple edge on
two parts, $K_4$ on four, and the prism $C_h\square K_2$ on $2h\ge6$. Map the three edges of $Q$ at each part
injectively to vertices of the part: any three vertices of a clique, and an attaching triple $A$ for $X$. Each edge
$PR$ of $Q$ becomes an edge of $G$ between the images of its ends. We call these edges \emph{cross edges}; every vertex
carries at most one.

\begin{theorem}\label{thm:construction}
The graph $G$ has $n$ vertices, $t(n,\alpha)+\ceil{3\alpha/2}$ edges, independence number $\alpha$, connectivity
$\kappa(G)\ge3$ and clique cover number $\theta(G)=\alpha+1$. In particular $V(G)$ has no partition into $\alpha$
cliques, and $s(G)=1$.
\end{theorem}

The proof uses six facts about the two blocks, each checked exhaustively by a short program (EXP-013,
\texttt{check\_lemmas.py}): for $X\in\{X_6,X_7\}$ and each attaching triple $A$,
\begin{enumerate}
\item[(F1)] $\alpha(X)=2$ and $\theta(X)=3$;
\item[(F2)] $X-A$ contains two non-adjacent vertices;
\item[(F3)] $\theta(X-S)+|S|/2\ge5/2$ for every nonempty $S\subseteq A$;
\item[(F4)] $X_6$ has $8=\e(K_3+K_3)+2$ edges and $X_7$ has $11=\e(K_3+K_4)+2$;
\item[(F5)] for every set $T$ of at most two vertices, every component of $X-T$ meets $A\setminus T$;
\item[(F6)] $X$ is $2$-connected.
\end{enumerate}

\begin{proof}
\emph{Edges.} The cliques of $G$ span $t(n,\alpha)$ minus the edges of the replaced pair, $X$ spans those edges plus
two (F4), and there are $3(\alpha-1)/2$ cross edges, which gives~\eqref{eq:count}.

\emph{Independence number.} An independent set meets each clique part in at most one vertex and $X$ in at most two
(F1), so $\alpha(G)\le\alpha$. For the converse, F2 gives non-adjacent $x,y\in X\setminus A$; neither carries a cross
edge, so neither has a neighbour outside $X$. The cross edges between clique parts form a graph of maximum degree at
most one whose vertex classes, the clique parts, have at least $3\ge2$ vertices; by Haxell's theorem \cite{Haxell}
there is a transversal of the clique parts with no cross edge inside it. With $x$ and $y$ it is an independent set of
size $\alpha$.

\emph{Connectivity.} Let $S$ be a set of at most two vertices. Every part keeps a vertex, and every clique part minus
$S$ is a nonempty clique. If $X-S$ is connected and $G-S$ were not, the parts would split into two nonempty sides with
every edge of $Q$ between them destroyed. But $Q$ has at least three such edges, and each vertex of $S$ destroys at most
its own cross edge. If $X-S$ is disconnected, then $S\subseteq V(X)$ by F6, so every other part and every cross edge
not at $X$ survives. Moreover $Q-X$ is connected: a cut vertex $v$ of a $3$-edge-connected cubic multigraph would leave
a side joined to the rest by at most one of the three edges at $v$. So the parts other than $X$ lie in one component.
By F5 each component of $X-S$ contains a vertex of $A\setminus S$, whose cross edge leads to a part other than
$X$, and so into the single component that holds those parts.

\emph{Clique cover number.} The $\alpha-2$ cliques and three cliques covering $X$ (F1) give $\theta(G)\le\alpha+1$. A
clique of $G$ that meets two parts is a single cross edge, since a third vertex would need a second cross edge at one
of its ends. Take a partition of $V(G)$ into cliques, let $M'$ be the set of cross edges it uses as cliques, and charge
each of them $1/2$ to each of its two parts. A part $P$ then carries
\[
c(P)=\theta\bigl(P-V(M')\bigr)+\tfrac12\,|P\cap V(M')|.
\]
A clique part carries at least $1$: either some vertex is left, needing one more clique, or all of its at least three
vertices are covered, giving at least $3/2$. For $X$, the covered vertices form a set $S\subseteq A$, and
$c(X)=\theta(X)=3$ when $S$ is empty and $c(X)\ge5/2$ otherwise, by F3. The partition has $\sum_Pc(P)\ge\alpha-2+5/2$
cliques, so at least $\alpha+1$.

\emph{Core parameter.} The parts of $G$ have independence numbers summing to $\alpha$, so $p(G)\ge\alpha-1$ by
Lemma~\ref{lem:p}, and $p(G)\ne\alpha$ because $\theta(G)>\alpha$.
\end{proof}

\begin{corollary}\label{cor:construction}
At each of the nine triples with $n\le24$ that it covers, $(10,3,3)$, $(11,3,3)$, $(16,5,3)$ to $(19,5,3)$ and
$(22,7,3)$ to $(24,7,3)$, the graph $G$ is extremal by Theorem~\ref{thm:values}, without any condition, and has no
partition into $\alpha$ cliques. If Didin's Theorem~9 holds, $G$ is extremal for every odd $\alpha\ge3$ and $3\alpha<n\le4\alpha-1$;
then case (a) of Proposition~\ref{prop:where} occurs at every triple where it is allowed, and the bound $n\le4\alpha-1$
there is sharp.
\end{corollary}

A program that builds $G$ with the triple edge ($\alpha=3$), $K_4$ ($\alpha=5$) or the prism $C_h\square K_2$
($\alpha\ge7$) as quotient and checks every property of
Theorem~\ref{thm:construction} exactly confirms it for every odd $\alpha\le11$ and every admissible order
(EXP-013, \nconstructed\ graphs). It also confirms that a quotient which is not $3$-edge-connected can fail: the
circular ladder on four parts, two double edges joined by two single edges, has a cut of two edges, and all
\nladder\ graphs built on it at $\alpha=5$, twenty random attachments at each order from $16$ to $19$, have the
conjectured number of edges and independence number $5$ but connectivity $2$ (EXP-013, \texttt{probe\_ladder.py}).
